\uuid{OO6t}
\titre{Solution générale/solution particulière}
\auteur{hanani ; mzari}
\organisation{exo7}
\licence{CC BY-NC-SA 4.0}
\datecreate{2020-01-20}
\niveau{L1}
\module{Analyse}
\chapitre{Equation différentielle}
\sousChapitre{Résolution d'équation différentielle du deuxième ordre}
\difficulte{3}
\theme{Equations du second ordre, solution générale, solution particulière}
\exo7qcmid{lille-2.equations-differentielles.032}

\contenu{
\begin{question}
\qclassification{225.02}
\qtags{motcle=solution générale/solution particulière}
On considère l'équation différentielle $(E)$ : $y''-2y'+2y=2\mathrm{e}^{x}\cos x$. Parmi les affirmations suivantes, cocher celles qui sont vraies :
\begin{answers}
\good{Les solutions de l'équation caractéristique sont $1\pm \mathrm{i}$.}
\bad{$(E)$ admet une solution sous la forme $\displaystyle y_0=a\mathrm{e}^{x}\cos x$ avec $a\in \Rr$.}
\good{La fonction $\displaystyle y_0=x\mathrm{e}^{x}\sin x$ est une solution de $(E)$.}
\bad{La solution générale de $(E)$ est : $\displaystyle y=\mathrm{e}^{x}\left(a\cos x+b\sin x\right)+2\mathrm{e}^{x}\cos x$, $a,b\in \Rr$.}
\end{answers}
\begin{explanations}
La solution générale de l'équation homogène est $Y=\mathrm{e}^{x}\left(a\cos x+b\sin x\right)$, $a$ et $b\in\Rr$, et $\displaystyle y_0=x\mathrm{e}^{x}\sin x$ est une solution particulière de $(E)$. Donc la solution générale de $(E)$ est : $\displaystyle y=Y+y_0=\mathrm{e}^{x}\left(a\cos x+b\sin x\right)+x\mathrm{e}^{x}\sin x$, $a,b\in \Rr$.
\end{explanations}
\end{question}
}
